% 原书分章 PDF 第 1—3 页：章首导言与第 28.1 节。
\chapter{靶向自助法}
\label{ch:targeted-bootstrap}
\chapterauthors{Jeremy Coyle 与 Mark J. van der Laan}
\chapteraffiliations{%
J. Coyle（通讯作者）：美国加利福尼亚大学伯克利分校生物统计学部，101 Haviland Hall, \#7358, Berkeley, CA 94720, USA；电子邮箱：\href{mailto:jrcoyle@berkeley.edu}{jrcoyle@berkeley.edu}。\\
M. J. van der Laan：美国加利福尼亚大学伯克利分校生物统计学部及统计学系，101 Haviland Hall, \#7358, Berkeley, CA 94720, USA；电子邮箱：\href{mailto:laan@berkeley.edu}{laan@berkeley.edu}。\\
本章 DOI：\url{https://doi.org/10.1007/978-3-319-65304-4_28}。}

\begin{sourcebox}
自助法（bootstrap）用于在各种情形下进行统计推断，包括构造置信区间和进行假设检验（Efron 和 Tibshirani，1993；Davison 和 Hinkley，1997）。在某些情形下，基于自助法的置信区间已被证明具有高阶准确性（higher-order accuracy），优于基于正态近似的 Wald 型区间（Hall，1988、1992；DiCiccio 和 Romano，1988）。因此，自助法已被广泛用于各种情境中的推断，尽管这些应用并非全都有理论支持。在靶向学习框架中，按通常方式应用自助法就会失效。我们将详细说明这一失效的原因，并提出靶向自助法（targeted bootstrap）作为解决方案，使其能够相合地估计抽样分布的前两阶矩。
\end{sourcebox}

假定我们希望采用靶向学习方法估计某个路径可微参数。典型流程是：先通过超级学习器（super learner）获得似然相关因子的初始估计，再使用靶向最小损失估计（targeted minimum loss-based estimation，TMLE）生成靶向估计。结合超级学习器和 TMLE，无需假定似然可以用简单参数模型来描述，就能对感兴趣的参数作出正确推断。由于 TMLE 是渐近线性估计量，我们可以采用正态近似，构造标准误基于影响曲线的 Wald 型置信区间。这些置信区间具有一阶准确性。人们可能希望改用非参数自助法，得到高阶正确的置信区间。然而，当 TMLE 的初始估计由超级学习器算法得到时，直接套用非参数自助法并无充分依据。我们将展示其表现不佳，因为超级学习器在非参数自助样本上的行为，与其在真实数据生成分布的样本上的行为不同，因而 TMLE 的行为也随之不同。因此，有必要发展适用于靶向学习情境的自助法。

我们将说明超级学习器出现这一行为差异的原因，并以靶向自助法给出解决方案。这是一种新的基于模型的自助法，从一个同时靶向于感兴趣参数及其影响曲线渐近方差的分布中抽样。在此过程中，我们将给出一种同时靶向于感兴趣参数及其渐近方差的 TMLE。直接使用渐近方差的靶向估计，可以用这个 TMLE 构造另一种 Wald 型置信区间；它也可用于构造渐近方差本身的置信区间。我们将展示靶向自助法置信区间的实际表现，并与 Wald 型置信区间以及其他自助方法生成的置信区间比较。

\section{问题陈述}

假定我们观测到随机变量 $O$ 的 $n$ 个独立同分布副本，其概率分布 $P_0$ 已知属于统计模型 $\mathcal M$。进一步假定，对于给定的参数映射 $\Psi:\mathcal M\to\mathbb R$，我们关心目标参数值 $\psi_0=\Psi(P_0)$ 的统计推断。考虑给定的估计量 $\hat\Psi:\mathcal M_{np}\to\mathbb R$，它将 $O_1,\ldots,O_n$ 的经验分布 $P_n$ 映射为 $\psi_0$ 的估计。假定估计量 $\psi_n=\hat\Psi(P_n)$ 在 $P_0$ 处渐近线性，其在 $P_0$ 处的影响曲线为 $O\mapsto D(P_0)(O)$，从而可写成
\[
 \psi_n-\psi_0=(P_n-P_0)D(P_0)+o_P(1/\sqrt n).
\]
此时，$\sqrt n(\psi_n-\psi_0)$ 依分布收敛到正态分布 $N(0,\Sigma^2(P_0))$，其中 $\Sigma^2:\mathcal M\to\mathbb R$ 定义为 $\Sigma^2(P)=PD(P)^2$，即影响曲线 $D(P)$ 在 $P$ 下的方差。

我们希望估计 $\psi_n$ 的置信区间。单侧置信区间由一个量 $\psi_{n,[\alpha]}$ 定义，使
$P_0(\psi_0<\psi_{n,[\alpha]})=\alpha$。双侧置信区间通常为等尾区间（equal-tailed interval），即两尾的误差相同：
\[
 P_0(\psi_0<\psi_{n,[\alpha/2]})
 =P_0(\psi_0>\psi_{n,[1-\alpha/2]})=\alpha/2.
\]
这样的区间可由一对单侧区间构造。利用前述渐近正态性，可以构造单侧 Wald 置信区间：定义方差估计量 $\hat\sigma_n^2=\Sigma^2(P_n)$，则端点为
\[
 \psi_{n,[\alpha],\mathrm{Wald}}
 =\psi_n-n^{-1/2}\hat\sigma_n\phi^{-1}(1-\alpha),
\]
其中 $\phi^{-1}(1-\alpha)$ 是标准正态分布的 $1-\alpha$ 分位数。类似地，双侧 Wald 置信区间为
\[
 \psi_n\pm\phi^{-1(1-\alpha/2)}n^{-1/2}\hat\sigma_n.
\]
这一方法忽略了余项 $o_P(1/\sqrt n)$，因此被称为一阶正确。

为给出一个具体的说明性例子，假定我们观测到 $O=(W,A,Y)\sim P_0$ 的 $n$ 个独立同分布观测，其中 $W$ 为基线协变量，$A\in\{0,1\}$ 为处理，$Y\in\{0,1\}$ 为结局。假定 $\mathcal M$ 是非参数模型，不对生成 $O$ 的分布作任何假设。目标参数 $\Psi:\mathcal M\to\mathbb R$ 是特定处理下的均值（treatment-specific mean），定义为
$\Psi(P)=E_PE_P(Y\mid A=1,W)$。令
$\bar Q(P)(W)=E_P(Y\mid A=1,W)$，$\bar g(P)(W)=E_P(A\mid W)$。
